\documentclass[11pt]{article}
\usepackage[T1]{fontenc}
\usepackage[margin=23mm]{geometry}
\usepackage{amsmath,amssymb,booktabs,hyperref}
\hypersetup{hidelinks}
\newcommand{\F}{\mathbb F_2}
\newcommand{\Z}{\mathbb Z}
\newcommand{\Q}{\mathbb Q}
\newcommand{\U}{(\mathbb R/\mathbb Z)}
\newcommand{\Sq}{\operatorname{Sq}}
\newcommand{\SSq}{\mathsf S}
\newcommand{\lift}[1]{\widetilde{#1}}
\newcommand{\ds}{d_s}
\newcommand{\OO}{\widehat O}
\newcommand{\PP}{\mathcal P}
\newcommand{\CA}{\mathrm{CA}}
\newcommand{\pathref}[1]{{\small\nolinkurl{#1}}}
\title{Crystalline spinless fermions:\\background, incoming gauges and upper extensions}
\author{Independent space-group calculation note}
\date{September 2026}
\begin{document}
\maketitle

\section{The physical convention and its exact representative}
Let $G$ be the full infinite affine space group and $V$ its real
three-dimensional linear representation. The supplied convention is
\begin{align}
 s_{\rm eff}&=s_{\rm crys}+w_1(V),\\
 \omega_{\rm eff}&=\omega_{\rm crys}+w_2(V)
       +w_1(V)\bigl(s_{\rm crys}+w_1(V)\bigr).
\end{align}
Thus crystalline spinless fermions, with
$s_{\rm crys}=\omega_{\rm crys}=0$, require
\begin{equation}\label{eq:pin}
 s=w_1(V),\qquad [\omega]=w_2(V)+w_1(V)^2.
\end{equation}
This is the Pin-minus background, and differs from the effective
$\omega=0$ convention of the crystalline spin-half calculation.
All group cohomology below is for $G$, including its translations.

There is a direct exact construction of a normalized cocycle for
\eqref{eq:pin}. For each point matrix $R$ set $J(R)=\det(R)R$. Since
$J=V\otimes\det V$ is oriented and
$w_2(J)=w_2(V)+w_1(V)^2$, its Spin double cover supplies the required
extension. The implementation uses the rational invariant metric
$\sum_R R^T R$ and the corresponding negative Clifford algebra.
An even rational element $q_R$ satisfying
$q_R e_i=J(R)(e_i)q_R$ is unique up to scalar. Fixing its first nonzero
coordinate to one gives
\begin{equation}
 q_Rq_S=\lambda(R,S)q_{RS},\qquad
 \omega(R,S)=\begin{cases}0,&\lambda(R,S)>0,\\1,&\lambda(R,S)<0.\end{cases}
\end{equation}
Positive normalization factors do not change this sign. In particular,
a mirror and a proper twofold rotation square to the central minus sign;
spatial inversion squares to plus one. The finite multiplication and
cocycle tables are pulled back through the actual affine projection.
They construct the background, not a substitute finite classification.

If the pulled-back $\omega$ is exact on $G$, the code constructs and
retains an actual binary one-cochain trivialization before choosing the
cohomologous representative $\omega=0$. Only after this certified change
of extension gauge are the zero-background shortcuts applicable.
The original point table and native trivialization remain evidence.

The convention comes from the supplied
\pathref{reference/legacy_stacking_context/06_examples.tex}.
The implementation is \pathref{gap/crystalline_background.g}; the
affine gauge reduction is in \pathref{gap/background_operations.g}.

\newpage
\section{The general-background cochain equations}
Write a three-dimensional state as $(n,B,C,\nu)$ with coefficients
\begin{equation}
 n\in C^1(G;\Z_s),\quad B\in C^2(G;\F),\quad
 C\in C^3(G;\F),\quad\nu\in C^4(G;\U_s).
\end{equation}
The complete tower satisfies
\begin{equation}\label{eq:tower}
 \ds n=0,\qquad dB=\omega\cup a+s\cup a^2,
 \qquad dC=Q(n,B),\qquad\ds\nu=\OO_5(n,B,C),
 \quad a=[n]_2.
\end{equation}
Binary lifts are taken after forming the whole binary expression.
Integer and rational products retain interval-cut signs and local-system
transport. Exact divisions precede modular reduction, and negative
carries use floor. For a possibly nonclosed cochain $z$, put
\begin{equation}
 \SSq^j(z)=z\cup_{|z|-j}z+z\cup_{|z|-j+1}dz.
\end{equation}
The second term is essential. With
$h=[\lfloor n/2\rfloor]_2$, $u=B+s\cup h$, and $W=\omega+s^2$,
one has $du=W\cup a$. The supplied native secondary representative is
\begin{align}\label{eq:Q}
 Q(n,B)={}&\SSq^2u+\omega\cup u+s\cup\SSq^1u
       +\zeta_{2,1}(W,a)\nonumber\\
 &+(W\cup_1W+s\cup W)\cup h\nonumber\\
 &+\bigl[(W\cup_1W)\cup_1s+s\cup(s\cup_1W)\bigr]\cup a,
\end{align}
where $\zeta_{2,1}(W,a)=\smile_{123134}(W,W,a,a)$ is the fixed
Cartan word. Replacing it by a cohomologous expression requires its
cochain dictionary.

The terminal formula is the complete supplied normalized expression
\begin{equation}\label{eq:O}
 \OO_5=\tfrac12[\SSq^2C+\omega\cup C]
 +\OO_5^{\gamma,\mathrm{ext}}(u)
 +\OO_5^{\gamma\psi}+\OO_5^\psi.
\end{equation}
It includes the fixed Alexander--Whitney/Eilenberg--Zilber completion,
integer carries, the full lifted Cartan bracket, and the coupling
\begin{equation}
 \frac1{16}\bigl[\lift W\cup\lift W+
      \lift W\cup_1\ds\lift W\bigr]\cup n.
\end{equation}
The integer $n$, not merely its parity, occurs here. The transcription
was checked against the unchanged supplied scalar evaluator, with
the corrected AW edge-transport normalization (v2). The large
fixed completion is specified by the supplied normalized manuscript and
the checked expression graph, rather than abbreviated as an unknown
group-dependent primitive.

For closed Majorana input $b$, the lower CA theory uses
\begin{equation}
 d c=P(b)=\Sq^2b+\omega\cup b+s\cup\Sq^1b,
 \qquad\OO_5^{\CA}(b,c)=\tfrac12[\SSq^2c+\omega\cup c]
                          +\OO_5^\gamma(b).
\end{equation}
Its supplied product has $b''=b+b'$,
$c''=c+c'+b\cup_1b'+s\cup(b\cup_2b')$, and
$\nu''=\nu+\nu'+U_4^{\CA}$. The full CA phase and product, including
all background terms, are mathematical inputs. Their closure identity
does not fix an arbitrary added closed product correction.

\newpage
\section{The full incoming integer gauge for unitary groups}
When $s=0$, $H^0(G;\Z)=\Z$ contributes an incoming integer gauge.
Its first image is $[\omega]$ in the Majorana layer, but quotienting
only that image misses the later CF and bosonic images. In the delivered
CA coordinates a full incoming generator is the flat lower state
\begin{equation}
 X=(\omega,0,F(\omega)).
\end{equation}
Its CF component uses a given unary normalization: for nonnegative
integer gauge $n_0$, the supplied MC and CF components are
$n_0\omega$ and $\lfloor n_0/2\rfloor(\omega\cup_1\omega)$.
Thus $n_0=1$ gives $C=0$. The bundle's
\pathref{source/lib/ss_fermion.gi}:66--76 supplies this input;
\pathref{note/stacking_note.tex}:94--120 gives the repo-to-CA dictionary,
which changes only the phase, keeping $(a,c)$ fixed. No old implementation
is loaded by the independent runtime. This statement uses nonnegative
inputs; inverses are taken in the complete CA group.
Here $F$ is an additive rational phase modulo one. All CA expressions
in this section use the fixed background $(\omega,s=0)$. Define the integer
Pontryagin cochain and binary Bockstein representative by
\begin{equation}
 \PP(\omega)=\lift\omega\cup\lift\omega
        +\lift\omega\cup_1d\lift\omega,
 \qquad r=\left[\frac{d\lift\omega}{2}\right]_2.
\end{equation}
The normalized universal local four-cochain $F$ is fixed by
\begin{align}
 dF&=\OO_5^{\CA}(\omega,0),\\
 2F+U_4^{\CA}(\omega,0;\omega,0)&=-\PP(\omega)/8\pmod1.
\end{align}
The second equation is the even-input calibration $R_0(A=2,b=0)$
from the supplied collaborator's \pathref{python/low_phases.py}, with
its Pontryagin sign. This is formula input, not an answer-table lookup.
It is stronger than closure alone.

The implementation determines a universal table of 64 sixteenth-valued
coefficients. All 1,024 local closed-$\omega$ five-simplices, all 64
square equations and all 23 degenerate four-simplex states are checked.
In this fixed-CF normalized local model, the remaining half-valued
closed ambiguity has dimension four: three exact directions and
$\omega^2/2$. The last class is $8X$. Thus $X\mapsto X+8X=9X$
does not change the cyclic subgroup generated by $X$; the exact
directions are gauges.
Closure or the square calibration alone is not claimed to select every
possible CF normalization.

Successive literal CA products give
\begin{align}\label{eq:powers}
 2X&=(0,r,-\PP/8),\\
 4X&=(0,0,-\PP/4+\tfrac12 r\cup_2r),\\
 8X&=(0,0,-\PP/2),\qquad 16X=0.
\end{align}
The extra fourth-power term has an explicit gauge:
\begin{equation}
 r\cup_2r=d\eta,\qquad
 \eta(0123)=(1-\omega_{012})\omega_{013}(1-\omega_{023})\pmod2.
\end{equation}
These are the leading coordinates of the canonical multiples
$X,2X,4X$, not unconditional formulas on every page's primitive source.
For example, if $2X$ already has exact CF component, its gauge-corrected
bosonic class determines the next image; one must not replace it by
$4X$. Each incoming page is defined on the preceding kernel, with
all required lower gauges retained.

The program first measures the full marked lower CA group $H$, then
forms $H/\langle X\rangle$ and recomputes its associated graded groups
from integer relation lattices. This handles all three incoming pages
at once. Ratios of the pre/post graded orders measure the successive
filtered images; their product equals the actual order of $X$.
For nontrivial $s$, $H^0(G;\Z_s)=0$ and this closed incoming subgroup
is absent. The relevant files are
\pathref{gap/pip_incoming_background.g} and
\pathref{gap/background_quotient.g}.

\newpage
\section{The surviving free integer lattice}
The free rank alone does not specify its primitive surviving lattice.
At nonzero $\omega$, neither the unitary zero lower tower nor the
zero-background infinite-dihedral construction can be assumed valid.
The supplied normalized formula instead gives, for a signed integral
cocycle $m$, the fixed-upper-choice identity
\begin{equation}\label{eq:shift}
 \OO_5(n+8m,B,C)-\OO_5(n,B,C)
        =\tfrac12 W\cup W\cup[m]_2\pmod1.
\end{equation}
At $n=16m$ the parity $a$ and half-parity $h$ both vanish. Thus
$B=C=0$ solves both binary sources. Applying \eqref{eq:shift} twice
gives the pointwise flat universal tower
\begin{equation}\label{eq:16}
 (n,B,C,\nu)=(16m,0,0,0).
\end{equation}
This proves $16H^1(G;\Z_s)$ survives for arbitrary background.
It does not prove that the analogous eightfold phase is zero.

For the actual finite-holonomy tangent background there is a stronger
\emph{existence} theorem. Since $[W]=w_2(V)$ and
$\rho_2p_1(V)=w_2(V)^2$, the class
\begin{equation}
 \gamma=p_1(V)\cup m\in H^5(G;\Z_s)
\end{equation}
is an integral lift of $[W]^2[m]_2$. It is torsion, because $p_1(V)$
is pulled back from the positive-degree integral cohomology of the
finite point group. Choose an integral twisted cocycle $\Gamma$
representing $\gamma$ and a binary four-cochain $\lambda$ with
\begin{equation}
 d\lambda=W\cup W\cup[m]_2+\rho_2\Gamma.
\end{equation}
There is a rational twisted four-cochain $Q_4$ satisfying
$\ds Q_4=\Gamma$. Hence
\begin{equation}
 \nu_8=\tfrac12(Q_4+\lift\lambda),\qquad
 \ds\nu_8=\tfrac12 W\cup W\cup[m]_2
              =\OO_5(8m,0,0)\pmod1.
\end{equation}
This proves $8H^1(G;\Z_s)$ survives. The phase primitive remains
necessary for an actual marked generator. We assert no fourfold
theorem, and no eightfold theorem for an unrelated arbitrary background.

For these three-dimensional point representations, a genuinely nonzero
tangent background with free integer classes has free rank one.
Indeed $H^1(G;\Q_s)$ is the determinant multiplicity in $V^*$.
If that multiplicity is at least two, a proper operation fixes two
independent vectors and therefore all three. An improper operation
acts by minus one on those two vectors and hence also on the third.
The point group is then contained in $\{I,-I\}$, whose required
extension class is zero. The higher-rank cases therefore use the
certified zero-background branch.

For rank one the surviving projection is $k\Z$, with $k\mid16$.
The implementation tests $k=1,2,4,8$, including every available torsion
shift until one survives, and retains \eqref{eq:16} as a universal
fallback. A torsion shift changes the $H^1$ class, not merely its gauge.
Each selected generator receives actual $B,C,\nu$ cochains. For the
tangent background the stronger theorem forces $k\leq8$; index sixteen
would require investigation. In rank $r$ the general index bound is
$8^r$, not eight.

\newpage
\section{The torsion square and extension bounds}
For nontrivial $s$, the canonical integer torsion representative is
$n=s$. The supplied first stacking product is
\begin{equation}\label{eq:E}
 E_2(n,n')=a\cup a'+s\cup(a\cup_1a'),\qquad
 a=[n]_2,\quad a'=[n']_2.
\end{equation}
The source manuscript calls its closed cup-product term proposed.
We treat \eqref{eq:E} as mathematical input; its normalization does
not follow from the obstruction identity alone. Since
$s\cup_1s=s$, $E_2(s,s)=0$. Before removing the integer gauge, the
square has integer component $2s$ and zero Majorana component.

To identify the gauge's leading component, use an interval height
$t$ equal to zero at the bottom and one at the top. On the cylinder set
\begin{equation}
 n_I=-\ds t,\quad [n_I]_2=dt,\quad
 B_I=\omega\cup t+s\cup(t\cup dt).
\end{equation}
Then $dB_I=\omega\cup dt+s\cup(dt\cup dt)$, exactly the primary
source in \eqref{eq:tower}. Its endpoints are $(0,0)$ and $(2s,\omega)$.
Removing this integer gauge from the square gives the rigorous leading
relation, conditional on \eqref{eq:E},
\begin{equation}\label{eq:leading}
 \text{Majorana projection of }2P=[\omega].
\end{equation}
The equation is in the final Majorana quotient. It does not fix the
CF and bosonic coordinates of $2P$. Changing a chosen full lift by a
lower state $L$ changes its square by $2L$ and leaves this projection
unchanged.

There is a useful first-product normalization check. The universal group
for two signed integer cocycles is
$U=\Z^2\rtimes C_2$, with simultaneous inversion. Its degree-two
binary cohomology has basis
$s^2,sa,sa',aa'$: the spectral-sequence dimension bound is four,
and the four reflection subgroups $(u,v,1)$ detect the independent
functions $1,u,v,uv$. Identity normalization removes the first three
closed corrections. An independent $\omega$ adds a factor
$K(\F,2)$ to $BU$; its only extra degree-two class is the
input-independent $\omega$, also removed by identity normalization.
Only the $aa'$ bit remains, up to a Majorana coboundary.

The supplied spinless mirror table states strong $\Z_{16}$. For the
corresponding signed $C_2$ background $\omega=s^2$, the independently
computed lower CA group is $\Z_8$, and \eqref{eq:leading} makes
$2P$ odd in that group. The alternative $aa'$ bit makes it even and
cannot give an element of order sixteen. Thus the separately supplied
physical normalization fixes this bit. Our own $\Z_{16}$ control uses
\eqref{eq:E} and is not a circular independent proof of that input.

For a general surviving $P$, let $H$ be the complete lower group.
An abstract extension by $\Z_2$ is determined by $[2P]\in H/2H$.
The code examines the \emph{entire} fiber compatible with
\eqref{eq:leading}, allowing the unspecified CF/bosonic carries.
If all representatives give the same Smith factors, the abstract group
is fixed without choosing those carries. A nonzero Majorana projection
alone does not imply uniqueness. Multiple resulting types are unresolved
possibilities, not a proof that every possible type is physically realized.
Even a unique abstract type does not supply a marked upper cochain
relation; its export keeps \pathref{fullUpperPhaseWitness=false}.

For unitary groups there is also a geometric bound on incoming $X$.
Here $V$ is oriented of rank three, so
$\PP(w_2(V))=\rho_4p_1(V)$. Since $p_1(V)$ is torsion, its image
under the coefficient map $k\mapsto-k/4$ into $\U$ is zero.
Together with the exact $r\cup_2r$ correction, this gives $4[X]=0$.
It does not prove incoming $d_4=0$: after allowed CF and MC-incoming
gauges, $2X$ could still be an order-two bosonic class. Vanishing of
$d_4$ in the result table must therefore come from the actual filtered
quotient certificates, not this order bound alone.

\newpage
\section{Calculation and validation boundaries}
Classification tests the primary source first, then the secondary
source modulo allowed Majorana choices, and finally the terminal
class modulo CF-primary and MC-secondary images. A nonzero raw
$Q$ or $\OO_5$ class can be canceled by those choices. Only an
obstruction remaining in the correct quotient kills the candidate.
When constructing a selected full tower, a changed Majorana choice
is followed by re-solving $C$, re-evaluating the complete terminal
formula, and solving its actual phase equation. Class vanishing alone
is not exported as a phase cochain.

The independent native resolution provides integer linear algebra,
comparison maps and the homotopy correction from native primitives
to the required bar equations. An exact native primary operation may
be cohomologous to the transferred bar operation without being the
same native vector. The closed-CF optimization retains its comparison
correction and the actual phase lift. Solved rational primitives have
no imposed sixteenth-denominator restriction.

The full nonzero-background terminal graph has 34,198 scalar nodes.
Exact restrictions to pointwise even integer $n$ and to the actual
representative $n=s$ have 7,786 and 16,008 nodes, respectively. The
even restriction retains the integer $n$ and is enabled only when all
coordinates of its integral $H^1$ combination are even, including
torsion coordinates. Integral bar transfer preserves that parity.
The sign restriction requires the actual sign cochain, not just an
equal cohomology class. Neither restriction permits dropping background
terms or lift carries. The v3 source \pathref{73e7bae9} also enables
the exact projected native cup-zero evaluation. These restrictions
change evaluation cost, not the formulas or the allowed towers.

The following evidence is already retained independently of final
space-group campaign counts:
\begin{itemize}
 \item Every point-group background for SG1--230 passed normalization,
 pair and triple cocycle checks; selected affine groups also passed
 native closure checks.
 \item General-background legal towers agree exactly with the unchanged
 supplied terminal oracle. The even and sign restrictions each passed
 64 further legal oracle/full-graph/restricted-graph comparisons and
 corresponding GAP controls; negative and large integer values are
 included in the even fixtures.
 \item The universal $X$ equations and ambiguity calculation are exact
 finite local checks. Actual incoming powers were tested on SG3, 16,
 75 and 143; full cyclic quotient controls include SG3 and 16.
 \item Actual finite $C_2$ with unitary $\omega=\chi^2$ has a lower
 $\Z_2$ killed by $X$. Signed $C_2$ with $\omega=s^2$ has measured
 lower relations $2D=0$, $2C=D$, $2B=C$; the permitted upper fiber
 has the unique abstract type $\Z_{16}$. No upper phase is inferred
 from that uniqueness.
\end{itemize}
Source hashes, fixtures and completed controls are in
\pathref{docs/validation_runs/}, especially the background,
incoming, quotient, finite-$C_2$ and O5-specialization records.
Detailed derivations are in
\pathref{docs/CRYSTALLINE_SPINLESS_BACKGROUND.md}.
This note records formulas and scope, not a final 230-group count or
performance benchmark. It does not claim a universal coherence proof
of the complete upper stacking law. Unknown upper carries remain
unknown unless an actual cochain formula or an exhaustive abstract
extension argument supplies the stated conclusion.

\end{document}
